\section{Theory Bridge and Claim Discipline}
\label{sec:ieee-theory}

The theorem surface of ORIUS is intentionally narrower than the rhetoric of a fully
closed universal program. The paper uses theory to justify the shape of the runtime
layer, not to imply that every domain has already inherited equal proof depth.

For readability, the main paper presents a narrow theorem spine rather than
inflating every supporting result into a separate headline. The full theorem
appendix records which results are main theorems, corollaries, definitions, and
lemmas; several are deliberately scoped: finite-horizon certificates, paired-trace
useful-work degradation, mandatory-release impossibility, two-state boundary lower
bounds, and adapter-semantics sensor necessity.

The bridge has three parts. First, if admissibility depends on the true state and the
observation channel can deviate materially from that state, then OASG is a genuine
safety event rather than a notation trick. Second, if lower reliability produces a
monotone inflation of the observation-consistent state set, then the admissible action
set must become smaller or remain unchanged. Third, if the runtime only releases actions
from that tightened set, or falls back when no such action exists, then the remaining
true-state violation risk is governed by the soundness of the uncertainty set and the
repair map rather than by the optimism of the nominal controller.

Temporal certificates and fallback add a fourth operational layer: T5 is a
finite-horizon certificate-validity theorem, T6 gives the confidence-aware
expiration bound, and T7\_Battery is the battery fallback instantiation.
Together they state that certificates are time-bounded safety objects and that
the battery runtime either holds safely, recovers by boundary-aware safe
landing, or fails closed.  The runtime enforces the resulting $\tau_t \ge 1$
check at each step.

\begin{table}[t]
\centering
\caption{Theory-to-evidence bridge used in the IEEE draft.}
\label{tab:ieee-theory-bridge}
\small
\begin{tabularx}{\columnwidth}{@{}p{0.26\columnwidth}X@{}}
\toprule
\textbf{Claim layer} & \textbf{What the paper is willing to claim}\\
\midrule
Structural & Degraded observation can create action-semantic failure even when the nominal policy is internally coherent.\\
Contractual & A typed runtime layer can convert reliability loss into tightened admissible action sets, repair, fallback, and certificate semantics.\\
Empirical & Battery is the primary validation domain; AV and healthcare provide bounded validation under explicit contracts.\\
Non-claim & The paper does not claim global minimax optimality, equal deployment readiness, or flat closure beyond the active three-domain evidence boundary.\\
\bottomrule
\end{tabularx}
\end{table}

This distinction is important because ORIUS aims to be scientifically ambitious without
being careless. The architecture is intentionally written as a universal safety layer
for Physical AI. The evidence, however, is tiered. That tiering is not a weakness in
the argument; it is the mechanism that keeps theorem language, benchmark language, and
runtime language aligned with tracked artifacts.

The practical consequence is straightforward. A domain does not enter the defended
evidence surface merely because it looks structurally compatible with the adapter
contract. It becomes defended only after the data surface, replay surface, repair
semantics, fallback semantics, certificate traces, and benchmark outputs satisfy the
claim-evidence criteria. That is
why the universal claim in this paper is architectural and bounded rather than flatly
equalized across domains already. The generated theorem-evidence matrix and
result cards are the status authority: T9/T10 are promoted only as scoped
mandatory-release and two-state lower-bound theorems, not as unrestricted
global impossibility or minimax-optimality claims.
